\documentclass[../../../include/open-logic-section]{subfiles}

\begin{document}

\olfileid{sth}{story}{urelements}
\olsection{بنسټيز غير سټ غړي ومنو که نۀ؟}

په راتلونکو څو څپرکو کښې به هڅه وکړو چې د سټ له
تجمعي-تکراري تصور څخه اصلونه راوباسو۔ خو تر نور وړاندې
تګ مخکې بايد د \emph{بنسټيزو غير سټ غړو} په باب څه نور ووايو۔

د \olref[approach]{sec} انځور يوازې دا اجازه راکوله چې له
زړو \emph{سټونو} نوي سټونه جوړ کړو۔ خو ښايي دا اجازه هم
وغواړو چې ځينې \emph{غير سټونه}، غواګانې، خوګان، د شګو
دانې يا بل هر څه، د سټونو !!{element}s شي۔ په دې حالت کښې
به له ځينو بنسټيزو غړو، \emph{بنسټيزو غير سټ غړو}، پيل
وکړو او بيا به ووايو چې په هر پړاو $S$ کښې هغه «ټول ممکن»
سټونه جوړوو چې له بنسټيزو غير سټ غړو يا په تر $S$ وړاندې
پړاوونو کښې له جوړ شويو سټونو څخه جوړ وي (په هر ترکيب
کښې)۔ پيدا شوے انځور به دې ته ډېر ورته وي:
\begin{center}
	\begin{tikzpicture}[scale=0.6]
	\tikzset{cut_here/.style={densely dotted}}
	\draw (-4,6) -- (-1, 0) -- (1,0) -- (4,6); 
	\draw[cut_here] (-5,8)--(-4,6);
	\draw[cut_here] (5,8)--(4,6);
	\draw(-1.5, 1)--(1.5, 1);
	\draw(-2, 2)--(2, 2);
	\draw(-2.5, 3)--(2.5,3);
	\draw(-3, 4)--(3, 4);
	\draw(-3.5, 5)--(3.5, 5);
	\draw(-4, 6)--(4, 6);
	\node[label] at (2, 0) {\small 0};
	\node[label] at (2.5, 1) {\small 1};
	\node[label] at (3, 2) {\small 2};
	\node[label] at (3.5, 3) {\small 3};
	\node[label] at (4, 4) {\small 4};
	\node[label] at (4.5, 5) {\small 5};
	\node[label] at (5, 6) {\small 6};
	\end{tikzpicture}
\end{center}
نو اوس يوه پرېکړه راته پکار ده: \emph{آيا بنسټيز غير سټ غړي ومنو؟}

له فلسفي پلوه په خپله تيوري کښې د بنسټيزو غير سټ غړو
راګډول معنا لري۔ اصلي وجه يې د سټ تيورۍ د \emph{کارولو وړ}
ګرځول دي۔ د دې ټکي د څرګندولو لپاره د \olref[sfr][siz][]{chap}
دا خبره راياده کړئ چې وايو دوه سټونه $A$ او~$B$ يو شان
اندازه لري، يعنې $\cardeq{A}{B}$، که او يوازې که تر منځ يې
يوه يو-پر-يو او پر-ټوله تابع وي۔ اوس که د پټي غواګانې او
د خوګانو په ځاے کښې خوګان دواړه سټونه جوړوي، نو د دې
ادعا لپاره يو سټ-تيوريکي چلند ورکولے شو چې «د غواګانو او
خوګانو شمېر يو شان دے»۔ خو که بنسټيز غير سټ غړي منع کړو،
نو غواګانې او خوګان سټونه \emph{نۀ} جوړوي او دا سټ-تيوريکي
چلند به نۀ وي۔ په رښتيا، د سټونو پخپله پرته به بل هيڅ شي
ته د سټ تيورۍ د کارولو نېغ وس نۀ لرو۔ (د بنسټيزو غير سټ
غړو د راګډولو د نورو وجهو لپاره \citealt[pp.~vi, 24,
50--1]{Potter2004} وګورئ۔)

خو له رياضيکي پلوه د بنسټيزو غير سټ غړو منل ډېر کم دي۔
يوه وجه يې دا ده چې بې له بنسټيزو غير سټ غړو د سټ تيورۍ
بيانول \emph{ډېر لږ} اسانه دي۔ خو کله کله له سټ تيورۍ د
بنسټيز غير سټ غړي د ايستلو لپاره لا په زړه پورې دليلونه
موندل کېږي:
\begin{quote}
  له دې باور سره سم چې سټ تيوري د رياضۍ بنسټ دے، بايد د
  ټولې رياضۍ بيان يوازې د سټونو په خبرو وکړے شو؛ نو زمونږ
  متغير بايد د غواګانو او خوګانو غوندې شيانو پر دامن نۀ ګرځي۔
	%But if $C$ is a cow, $\{C\}$ is a set, but not a legitimate mathematical object. 
  \citep[p.~8]{Kunen1980}
\end{quote}
نو د کارېدنې پام د بنسټيزو غير سټ غړو \emph{راګډول} وړانديز
کوي؛ د بنسټ لپاره تحويلي موخه (خالصې سټ تيورۍ ته د رياضۍ
راکمولو موخه) د هغو \emph{ايستل} وړانديز کولے شي۔ لږه
سستي هم د بنسټيزو غير سټ غړو د ايستلو لوري ته اشاره کوي۔

مونږ به تر ټولو اسانه لاره ونيسو۔ خو يوه وجه يې تعليمي هم
ده۔ موخه مو دا ده چې د سټ تيورۍ له هغو بنسټونو سره مو اشنا
کړو چې د سټ تيورۍ د فلسفې د پيل لپاره يې اړتيا لرئ۔ او د
دغې فلسفې زياتره ليکنې هغه سټ تيورۍ څېړي چې \emph{بې له}
بنسټيزو غير سټ غړو بيان شوې دي۔ نو دا کتاب به ښايي ډېر
ګټور وي که له دغو ليکنو سره نژدې پاتې شي۔

\end{document}
